\chapter{}

Havia mais uma razão para que eu decidisse retomar minha vida ao lado do Paulo.
Durante toda a doença do meu pai, ele pacientemente se pusera em segundo plano por anos, até que o Comandante se fosse desta vida.
Agora chegara a vez dele, Paulo, de viver situação semelhante.
E enfrentando condições muito piores.
Entre os filhos de José Ignácio, ele sempre fora o que mais se atritara com o pai e as discussões entre os dois, que umas poucas vezes me fora dado presenciar, eram de tal violência que chocavam até mesmo a mim, que não era propriamente neófita em conflitos familiares.
No entanto, mortos os irmãos, só restavam ele e a irmã Marina para assumir os cuidados com o velho, em acelerado processo de decadência física e mental.
Marina, de saída, eximiu-se da tarefa e com razão: ela e o marido, Omar, haviam cuidado por um longo tempo de Dona Yolanda.
José Ignácio seria missão do Paulo.
E eu fui testemunha de que ele a cumpriu até o fim, irrepreensivelmente, sem nenhum rancor, com absoluto zelo pela dignidade do pai.
Deve ter doído em José Ignácio descobrir que não fez por merecer o filho que teve.
Aliás, nenhum dos filhos que teve.
Mas, como ensinam os jesuítas, nunca sabemos como Deus trabalha no interior dos homens.
O fato é que era por Paulo que meu sogro chamava incessantemente nos seus últimos dias e só se acalmava quando o via diante de si.

Quando fui a Cananéia pela primeira vez, a passeio, comuniquei ao Paulo que moraria lá, sim, com a condição de que fosse de frente para o mar.
Com paciência evangélica, ele conseguiu comprar o único terreno disponível à margem do canal e que lá estava há anos, cercado por um muro meio desabado de pedras seculares, porque faltava coragem aos interessados para negociar com o proprietário, o Jornalista Antonio Santiago, um memorialista octogenário e impiedoso, capaz de massacrar um cristão por horas a fio com intermináveis digressões sobre acontecimentos e personagens do passado da cidade, do estado e do país que ele, como jornalista, historiador e aposentado da \textit{``Light and Power Co.''}, estava em melhores condições de conhecer que ``muitos autores mequetrefes'' por aí, que nada sabem da nossa História.
E, como se não bastasse, fazendo absoluta questão de documentar suas lembranças com toneladas de recortes, fotos e cartas que amareleciam e mofavam em pilhas espalhadas por todos os cantos do sombrio casarão da praça, onde fomos encontrá-lo.


Por meses, Paulo pagaria todos os seus pecados dividindo seu tempo entre José Ignácio e Antônio Santiago.
Como recompensa, acredito, Deus concedeu que o Jornalista, após longo e insistente assédio, finalmente se calasse tempo suficiente para ouvir a proposta de compra do terreno.
Nessa altura, ele se afeiçoara tanto ao Paulo que dizia, a torto e a direito, que esse era o filho que ele nunca tivera.
Aceitou a proposta e ainda nos deu de presente o portão de ferro centenário instalado em nosso jardim e que por anos ele guardara, na intenção de colocá-lo à entrada das terras da família, quando a interminável contenda jurídica que lhe tomara a vida e os recursos fossem devolvidas a ele e aos seus irmãos.

A compra do terreno concretizou-se, por fim, às vésperas de eu deixar Teresina e vir morar em Cananéia.
E, a partir de então, eu começaria a pensar dia e noite na casa que construiríamos sobre ele.
Na verdade, reconstruiríamos, porque o local era tombado e teríamos que ser fiéis à arquitetura do prédio que originalmente ocupara aquele espaço.
Quando voltei a Cananéia, dessa vez para ficar, as obras começaram.

Eu tinha a casa nítida dentro da minha cabeça.
A vida inteira eu desejara construir uma e sabia agora, exatamente, como a queria.
Ela teria que ser a soma de todas as nossas experiências e lembranças, comportando impressões que nos ficaram da infância, dos lugares por onde andamos, imaginamos ou sonhamos.
Essa casa seria principalmente dedicada ao Paulo.
A família dele é mais longeva que a minha e me incomodava a idéia de que quando eu me fosse, ele ficasse sem um lugar bonito para viver.
Porque se houve alguma coisa com que nunca me conformei foi com o jeito esquisito de morar dos meus sogros: ambas as casas eram um amontoado triste de móveis, papéis e livros antigos o que, para uma descendente dos Filpi e dos Vitta, sugeria desencanto e baixa auto-estima.
Eram casas de velhos.
Não, eu queria luz, sol, perfume de flores, as cortinas brancas, a simplicidade franciscana e o cheiro de limpeza das casas da minha infância.
Porém, guardávamos muita coisa da gente do Paulo, objetos, mobílias e enfeites e ainda queríamos que também nossos filhos sentissem a casa como sua, reconhecendo nela sua própria história.
Parecia muito para um espaço tão pequeno.

\begin{figure}
\centering
\includegraphics[width=0.8\linewidth]{42/casa-cananéia.png}
\caption{A casa de Cananéia.}
\end{figure}

\begin{figure}
\centering
\includegraphics[width=0.8\linewidth]{42/o-canal.png}
\caption{O canal defronte à casa.}
\end{figure}

\begin{figure}
\centering
\includegraphics[width=0.8\linewidth]{42/paisagem.png}
\caption{A paisagem vista através das nossas janelas.}
\end{figure}

Quando a casa ficou pronta, eu a fui preenchendo com todas as coisas que precisavam estar nela.
Digo sempre que foi uma construção abençoada, porque tudo foi se encaixando nos seus lugares, como se sempre tivesse estado ali.
E a atmosfera que hoje a casa irradia é exatamente a que eu imaginei.
Apenas o objetivo de fazer dela um lugar onde Paulo se sentisse bem, tenho dúvidas de que venha um dia a se cumprir.
Ele é compulsivamente inquieto e instável para isso
Eu, contudo, se me fosse dado escolher, queria ficar aqui diante dessa janela, olhando os barcos coloridos balançando sonolentos nas mansas águas do canal.
Vendo os papagaios e guarás voando em formação lá para os lados do morro de São João em busca de pouso e a lua cheia, no começo da noite, nascer por detrás da Ilha Comprida e aos poucos erguer-se no céu, abrindo sobre as águas seu esplêndido manto de puro ouro.
De novo, me vêm à lembrança os parques de diversão da minha infância.
Se meu tempo está se esgotando, se o dono do brinquedo está prestes a virar na minha direção seu olhar inexorável e indiferente, que os dias que me restam possam se arrastar com a preguiça dessas ondas quebrando molemente nas pedras.
Porque eu gosto da Vida e volto a dizer que só me aflige a consciência de ter recebido dela mais do que dei.


Mamãe morreu há dois anos.
Um pouco antes disso esteve aqui, por duas vezes.
Gostava de caminhar comigo à tardinha, ao longo do passeio que margeia o canal defronte à nossa casa.
Olhando tudo à sua volta, dizia e repetia, sinceramente admirada:

\textit{``-- Filha, aqui você tem tudo o que realmente importa nessa vida.
Esse lugar me faz sentir a presença de Deus.'' }

Percebi que ela chegara ao limiar do tempo em que começamos a buscar a essência das coisas.
Cananéia, pequena, pobre e acanhada, jamais a teria atraído há apenas um par de anos atrás.

Uma coisa notável que me aconteceu, desde quando comecei a registrar essas memórias, foi descobrir que algumas pessoas irrompem nelas com uma força insuspeitada.
Minha mãe, mais do que todos.
Sempre achei que fui mais influenciada por meu pai do que por ela ao longo da vida.
Mas, tive que reconsiderar em parte esse juízo quando mamãe fez oitenta anos e meus irmãos pediram que eu escrevesse algumas palavras para serem lidas na festa em homenagem a ela, à qual não fui, porque na ocasião não podia me afastar do Colégio.

Quando me sentei para redigir o pequeno discurso, quis fugir ao vazio laudatório característico dessas homenagens.
Decidi que usaria a oportunidade para oferecer aos meus irmãos e, principalmente a ela própria, uma visão justa e honesta da Lectícia que eu, como filha mais velha, provavelmente conheci melhor que os outros.
Lá pelas tantas, saiu-me num impulso a frase: ``muito cedo sagrou-se sacerdotisa na religião da Vida''.
Foi aí que me dei conta de que o legado bendito desse devoto amor à Vida, em todas as suas manifestações, me veio dela.
O mais importante de todos que recebi.
Tenho em mim, é verdade, muito do meu pai.
Mas do realismo que tangenciava o pessimismo e da tendência à depressão, salvou-me a fé indestrutível da minha mãe em que sempre é possível tornar tudo melhor, mais alegre e mais bonito.

\begin{figure}
\centering
\includegraphics[width=0.8\linewidth]{42/mamãe+eu.png}
\caption{Mamãe e eu fotografadas pela Paula.}
\end{figure}

Em 10 de dezembro de 2008, ela morreu.
Não consegui chorar no dia.
Ao contrário, quase me alegrei.
Ela não merecia suportar as limitações de um organismo minado pela sucessão de tromboses que sofreu, a confusão mental que a humilhava e, principalmente, a perda drástica daquela aparência vibrante de que tanto se orgulhava.
Odiava envelhecer.

Arrumei-a no caixão, deixando-a bonita como ela me fez prometer que faria.
E o que senti enquanto a contemplava, rejuvenescida, serena e adormecida, foi uma profunda cumplicidade.

 
Nos dias em que o vento arrepia o mar e ele se atira bravo sobre a areia da praia, lançando no ar nuvens de borrifos em desafio aos imponentes contrafortes da Serra do Mar, dou razão a mamãe.
É muito fácil sentir aqui a presença de Deus.
E, diante da grandiosidade do que ele criou, dar-se conta da pequenez do homem.
Isso traz uma enorme paz.
Somos tão efêmeros, tão miúdos, que toda ansiedade em torno de realizações pessoais e legados para a posteridade parece risível.
A idade e a experiência me convenceram de que a Vida é um fantástico plano para cuja execução fomos convidados a colaborar, quase todos em humilde anonimato e apenas alguns, muito poucos, guindados ocasionalmente à fama, merecida ou não, e isso, afinal, importa bem pouco.
 O que conta é que todos, cada um a seu modo, temos um papel a desempenhar nessa construção.
Claro que há os que se recusam a cooperar, insatisfeitos com seu quinhão, como há também os tolos que se iludem, lá no alto, pensando que tudo acontece para seu próprio triunfo e regozijo.
Como há ainda os prepotentes que, ao longo da História, vez ou outra, tentam ditar os rumos desse projeto que começou muito antes deles e vai se estender para muito adiante das suas fugazes existências.
E que, na desastrada tentativa, interrompem, atrasam e desvirtuam uma obra cujo sentido e beleza, estupidamente, nem ao menos se esforçam por tentar vislumbrar.


Conta uma lenda antiga que um tapeceiro, já muito velho, tecia, dia após dia, um enorme tapete que parecia não ter fim.
Um vizinho, acompanhando o trabalho dedicado e primoroso do ancião, perguntou-lhe, certa vez:

\textit{``-- Meu bom homem, não receia que a morte o alcance antes que consiga terminar a sua obra?''.}

\textit{``-- Não''}, respondeu mansamente o velho,\textit{ ``pois não fui eu que a comecei''}.

Gostei dessa história quando a li, porque fala do que me ensinou a minha gente, os Filpi, os Vitta e os Credidio.
Ser grato e feliz por ter vindo ao mundo, tendo em mente que o importante aqui não é você, mas essa maravilhosa obra que é a Vida e que devemos nos empenhar, geração após geração, em levar adiante.
Ajudar, facilitar, fazer sempre o melhor que puder, indo além de si mesmo quantas vezes for necessário, não esquecendo, entretanto, que nem a obra e nem a glória nos pertencem.
Dando-se até o direito de sentir uma leve tristeza por sair antes do fim, mas encontrando consolo na certeza de ter feito tudo de que foi capaz para engrandecê-la.

Muitas vezes, ouvindo os velhos da nossa família, percebi quanto ansiavam por poder deixar às gerações seguintes a receita do viver.
 Não conheci quem a tivesse e os mais sábios garantem que não há, pois para que houvesse uma, teríamos que saber primeiro a quê a Vida se destina.
Parece que temos de aceitar que faz parte da experiência levar conosco a inquietação de nunca alcançar-lhe o sentido.
Contudo, uma santa, certamente inspirada por Deus, talvez tenha encontrado uma fresta iluminada nesse escuro mistério. Ela ensinou:

\vfill

{
    \large
    \centering
    \itshape
    Muitas vezes as pessoas são egocêntricas,\\
    ilógicas e insensatas.

    Perdoe-as assim mesmo.

    Se você é gentil, as pessoas podem acusá-lo de\\
    egoísta, interesseiro.

    Seja gentil assim mesmo.

    Se você é vencedor, terá alguns falsos\\
    amigos e alguns inimigos verdadeiros.

    Vença assim mesmo.

    Se você é honesto e franco, as pessoas\\
    podem enganá-lo.

    Seja honesto e franco assim mesmo.

    O que você levou anos para construir,\\
    alguém pode destruir de uma hora para outra.

    Construa assim mesmo.

    O bem que você faz hoje,\\
    pode ser esquecido amanhã.

    Faça o bem assim mesmo.

    Dê ao mundo o melhor de você,\\
    mas, isso pode nunca ser o bastante.

    Dê o melhor de você assim mesmo.

    Veja você que, no final das contas,\\
    é entre você e Deus.\\
    Nunca foi entre você e as outras pessoas

}

\vfill

\begin{flushleft}
    \textbf{-- Madre Teresa de Calcutá}
\end{flushleft}

\vfill

\begin{flushright}
    Cananéia, 14 de novembro de 2010.
\end{flushright}